% GENERATED FILE. Edit canonical lessons, then run npm run generate:books.
% canonical-source-hash: fnv1a64:a168f1061de7376d
% canonical-lessons: MW-C193-rupiyo, MW-C193-sikko, MW-C193-not, MW-C193-nagad, MW-C193-chukavno

\chapter{Rupees and Coins}
\label{ch:mw-rupees-and-coins}
\hlchaptermodality{193}

\begin{chapteropening}
\textbf{By the end of this chapter:} \emph{I can count out rupees and coins and pay.}

Complete the last lesson of chapter 193: I can count out rupees and coins and pay.
\end{chapteropening}

\section[\texorpdfstring{\mw{रुपियो} (rupiyo)}{रुपियो (rupiyo)}]{\mw{रुपियो} (rupiyo) — a rupee}

\label{lesson:MW-C193-rupiyo}



\begin{quote}
Before the new one: say the Marwadi for to print, then the Marwadi for to mend.
\end{quote}

\subsection*{You'll want to know: \mw{रुपियो}}
\textbf{\mw{रुपियो}} — \emph{rupiyo} — \textquotedblleft{}a rupee\textquotedblright{}.

\begin{quote}
\mw{म्हारी} \mw{जेब} \mw{मांय} \mw{एक} \mw{रुपियो} \mw{है।}
\end{quote}

\begin{quote}
\emph{mhārī jeb mā̃y ek rupiyo hai.}
\end{quote}

\begin{quote}
I have one rupee in my pocket.
\end{quote}

A masculine noun.

\begin{cousinweb}
\textbf{\mw{रुपियो}} (\emph{rupiyo}) is inherited from Sanskrit \emph{rūpya}, \textquotedblleft{}bearing a stamp; silver\textquotedblright{}; CDIAL lists Marwari \emph{rūpo}, \textquotedblleft{}silver\textquotedblright{}, under it, and sets up a second form, \emph{*rūpiya-}, which the coin's name probably continues (CDIAL 10805).
\end{cousinweb}

\begin{grammarlens}[title={What you've built}]
One more word for reading notices and messages.
\end{grammarlens}

\subsection*{Your turn}
\begin{itemize}
\raggedright
  \item \emph{Say it:} \emph{rupiyo}
  \item \emph{Say it:} \emph{rupiyo}, once more
  \item \emph{Say it:} \emph{mhārī jeb mā̃y ek rupiyo hai.} — I have one rupee in my pocket.
  \item \emph{Recall:} say the Marwadi for to print, then the Marwadi for to mend, then say \emph{rupiyo} again
\end{itemize}

\begin{tcolorbox}[breakable,colback=teal!4,colframe=teal!35!black,title={Before you move on}]
What does \mw{रुपियो} mean? (\textquotedblleft{}A rupee\textquotedblright{}.) Say it once more.
\end{tcolorbox}
\section[\texorpdfstring{\mw{सिक्को} (sikko)}{सिक्को (sikko)}]{\mw{सिक्को} (sikko) — a coin}

\label{lesson:MW-C193-sikko}



\begin{quote}
Before the new one: say the Marwadi for to mend, then the Marwadi for a rupee.
\end{quote}

\subsection*{You'll want to know: \mw{सिक्को}}
\textbf{\mw{सिक्को}} — \emph{sikko} — \textquotedblleft{}a coin\textquotedblright{}.

\begin{quote}
\mw{ये} \mw{सिक्को} \mw{भारी} \mw{है।}
\end{quote}

\begin{quote}
\emph{ye sikko bhārī hai.}
\end{quote}

\begin{quote}
This coin is heavy.
\end{quote}

A masculine noun.

\begin{cousinweb}
\textbf{\mw{सिक्को}} (\emph{sikko}) is usually traced, through Persian, to Arabic \emph{sikka}, \textquotedblleft{}a die for stamping coins\textquotedblright{}: the coin is named for its stamp.
\end{cousinweb}

\begin{grammarlens}[title={What you've built}]
One more word for reading notices and messages.
\end{grammarlens}

\subsection*{Your turn}
\begin{itemize}
\raggedright
  \item \emph{Say it:} \emph{sikko}
  \item \emph{Say it:} \emph{sikko}, once more
  \item \emph{Say it:} \emph{ye sikko bhārī hai.} — This coin is heavy.
  \item \emph{Recall:} say the Marwadi for to mend, then the Marwadi for a rupee, then say \emph{sikko} again
\end{itemize}

\begin{tcolorbox}[breakable,colback=teal!4,colframe=teal!35!black,title={Before you move on}]
What does \mw{सिक्को} mean? (\textquotedblleft{}A coin\textquotedblright{}.) Say it once more.
\end{tcolorbox}
\section[\texorpdfstring{\mw{नोट} (noṭ)}{नोट (noṭ)}]{\mw{नोट} (noṭ) — a banknote}

\label{lesson:MW-C193-not}



\begin{quote}
Before the new one: say the Marwadi for a rupee, then the Marwadi for a coin.
\end{quote}

\subsection*{You'll want to know: \mw{नोट}}
\textbf{\mw{नोट}} — \emph{noṭ} — \textquotedblleft{}a banknote\textquotedblright{}.

\begin{quote}
\mw{म्हारो} \mw{नोट} \mw{नुवो} \mw{है।}
\end{quote}

\begin{quote}
\emph{mhāro noṭ nuvo hai.}
\end{quote}

\begin{quote}
My banknote is new.
\end{quote}

A masculine noun.

\begin{cousinweb}
\textbf{\mw{नोट}} (\emph{noṭ}) is borrowed from English \emph{note}, as in a bank note.
\end{cousinweb}

\begin{grammarlens}[title={What you've built}]
One more word for reading notices and messages.
\end{grammarlens}

\subsection*{Your turn}
\begin{itemize}
\raggedright
  \item \emph{Say it:} \emph{noṭ}
  \item \emph{Say it:} \emph{noṭ}, once more
  \item \emph{Say it:} \emph{mhāro noṭ nuvo hai.} — My banknote is new.
  \item \emph{Recall:} say the Marwadi for a rupee, then the Marwadi for a coin, then say \emph{noṭ} again
\end{itemize}

\begin{tcolorbox}[breakable,colback=teal!4,colframe=teal!35!black,title={Before you move on}]
What does \mw{नोट} mean? (\textquotedblleft{}A banknote\textquotedblright{}.) Say it once more.
\end{tcolorbox}
\section[\texorpdfstring{\mw{नगद} (nagad)}{नगद (nagad)}]{\mw{नगद} (nagad) — in cash}

\label{lesson:MW-C193-nagad}



\begin{quote}
Before the new one: say the Marwadi for a coin, then the Marwadi for a banknote.
\end{quote}

\subsection*{You'll want to know: \mw{नगद}}
\textbf{\mw{नगद}} — \emph{nagad} — \textquotedblleft{}in cash\textquotedblright{}.

\begin{quote}
\mw{हूं} \mw{नगद} \mw{देऊं} \mw{हूं।}
\end{quote}

\begin{quote}
\emph{hū̃ nagad deū̃ hū̃.}
\end{quote}

\begin{quote}
I pay in cash.
\end{quote}

\begin{cousinweb}
\textbf{\mw{नगद}} (\emph{nagad}) is usually traced, through Persian, to Arabic \emph{naqd}, \textquotedblleft{}ready money\textquotedblright{}.
\end{cousinweb}

\begin{grammarlens}[title={What you've built}]
One more word for reading notices and messages.
\end{grammarlens}

\subsection*{Your turn}
\begin{itemize}
\raggedright
  \item \emph{Say it:} \emph{nagad}
  \item \emph{Say it:} \emph{nagad}, once more
  \item \emph{Say it:} \emph{hū̃ nagad deū̃ hū̃.} — I pay in cash.
  \item \emph{Recall:} say the Marwadi for a coin, then the Marwadi for a banknote, then say \emph{nagad} again
\end{itemize}

\begin{tcolorbox}[breakable,colback=teal!4,colframe=teal!35!black,title={Before you move on}]
What does \mw{नगद} mean? (\textquotedblleft{}In cash\textquotedblright{}.) Say it once more.
\end{tcolorbox}
\section[\texorpdfstring{\mw{चुकावणो} (chukāvṇo)}{चुकावणो (chukāvṇo)}]{\mw{चुकावणो} (chukāvṇo) — to pay; to settle (a bill)}

\label{lesson:MW-C193-chukavno}



\begin{quote}
Before the new one: say the Marwadi for a banknote, then the Marwadi for in cash.
\end{quote}

\subsection*{You'll want to know: \mw{चुकावणो}}
\textbf{\mw{चुकावणो}} — \emph{chukāvṇo} — \textquotedblleft{}to pay; to settle (a bill)\textquotedblright{}.

\begin{quote}
\mw{आप} \mw{पैसा} \mw{अबार} \mw{चुकावो।}
\end{quote}

\begin{quote}
\emph{āp paisā abār chukāvo.}
\end{quote}

\begin{quote}
Please pay now.
\end{quote}

Past (I …): \mw{म्हैं} \mw{पैसा} \mw{चुकाया}, \emph{mhaĩ paisā chukāyā}.

\begin{cousinweb}
\textbf{\mw{चुकावणो}} (\emph{chukāvṇo}) is the causative of a verb inherited from a reconstructed \emph{*cukk}, \textquotedblleft{}fall short of, stop\textquotedblright{}: CDIAL lists Marwari \emph{cūkṇo}, \textquotedblleft{}to be completed\textquotedblright{}, and Hindi \emph{cukānā}, \textquotedblleft{}to settle\textquotedblright{} (CDIAL 4848). To pay a bill is to make it complete.
\end{cousinweb}

\begin{grammarlens}[title={What you've built}]
One more word for reading notices and messages.
\end{grammarlens}

\subsection*{Your turn}
\begin{itemize}
\raggedright
  \item \emph{Say it:} \emph{chukāvṇo}
  \item \emph{Say it:} \emph{chukāvṇo}, once more
  \item \emph{Say it:} \emph{āp paisā abār chukāvo.} — Please pay now.
  \item \emph{Recall:} say the Marwadi for a banknote, then the Marwadi for in cash, then say \emph{chukāvṇo} again
\end{itemize}

\begin{tcolorbox}[breakable,colback=teal!4,colframe=teal!35!black,title={Before you move on}]
What does \mw{चुकावणो} mean? (\textquotedblleft{}To pay; to settle (a bill)\textquotedblright{}.) Say it once more.
\end{tcolorbox}
